\subsection{可忽略的正函数}

定义 1. --- 给定一个测度 \(\mu\) 于局部紧空间 \(X\) 上，数值函数 \(f \geqslant 0\) （有限或否，定义在 \(X\) 上）称为关于测度 \(\mu\) 可忽略的，如果 \(|\mu|^*(f) = 0\) 。

这时我们也称 f 是 \(\mu\) -可忽略的，或在不致混淆时简称可忽略的。

命题 1. --- 若 \(f\) 是一个可忽略的函数 \(\geqslant 0\) ，则每个数值函数 \(g\) （满足 \(0 \leqslant g \leqslant \alpha f\) ， \(\alpha\) 为标量 \(> 0\) ）都是可忽略的。

事实上，\(0 \leqslant |\mu|^*(g) \leqslant \alpha |\mu|^*(f) = 0\)

命题 2. --- 序列 \((f_{n})\) （诸函数为可忽略的 \(\geqslant0\) ）的和与上包络都是可忽略的。

事实上，\(|\mu|^*(\sum_n f_n) \leqslant \sum_n |\mu|^*(f_n) = 0\) （§1, No.~3, 命题 13），且 \(\sup_n f_n \leqslant \sum_n f_n\) 。

命题 3. --- 下半连续函数 \(f \geqslant 0\) （定义在 \(X\) 上）可忽略，必须且只须 \(f\) 在 \(\mu\) 的支集上为零。

若 \(|\mu|^*(f) = 0\) ，则 \(|\mu|(g) = 0\) 对每个函数 \(g \in \mathcal{H}_+\) （满足 \(g \leqslant f\) ）成立；由此（Ch. III, §2, No.~3, 命题 9）推出 \(g\) 在 \(\mu\) 的支集 S 上为零；由于 \(f\) 是诸函数 \(g \in \mathcal{H}_+\) （满足 \(g \leqslant f\) ）的上包络（§1, No.~1, 引理），故在 S 上 \(f(x) = 0\) 。反之，若在 S 上 \(f(x) = 0\) ，则在 S 上 \(g(x) = 0\) 对每个函数 \(g \in \mathcal{H}_+\) （满足 \(g \leqslant f\) ）成立，因此（Ch. III, §2, No.~3, 命题 8）\(|\mu|(g) = 0\) ，按定义这蕴含 \(|\mu|^*(f) = 0\) 。

